\section{Results}
\label{sec:results}

\subsection{Water level}
\label{sec:res_wl}

Table~\ref{tab:wl} reports water level skill at the three interior stations for
every ensemble member over the post-spinup window. Raw error is dominated by a
positive bias at BocaNord and AltaVilaEst, between 0.09 and 0.18~m, and is
small at BocaSud, between $-0.004$ and $+0.018$~m. Removing the respective
means leaves an anomaly RMSE of 0.021 to 0.045~m at every station in every
member, with correlations from 0.79 to 0.93 and modelled-to-observed standard
deviation ratios from 1.06 to 1.25.

The spread across configurations is small. Anomaly RMSE varies by 0.005~m or
less between members at any given station, and the ordering of members is not
consistent from one station to another. Variable roughness lowers the
AltaVilaEst bias from 0.117 to 0.091~m relative to its own counterpart and
raises the BocaNord bias from 0.172 to 0.177~m.

\begin{table}[H]
  \centering
  \caption{Water level skill by member and station, 2--10 July 2025.
  $\mathrm{RMSE}^2 = \mathrm{bias}^2 + \mathrm{RMSE}_{\mathrm{anom}}^2$. All
  values in metres except the dimensionless ratios.}
  \label{tab:wl}
  \begin{tabular}{llrrrrr}
    \toprule
    Station & Member & RMSE & Bias & RMSE$_{\mathrm{anom}}$ & $r$ & $\sigma_m/\sigma_o$ \\
    \midrule
    BocaNord & no waves    & 0.150 & $+0.146$ & 0.034 & 0.86 & 1.09 \\
             & waves       & 0.170 & $+0.168$ & 0.030 & 0.88 & 1.09 \\
             & + roughness & 0.172 & $+0.170$ & 0.030 & 0.89 & 1.10 \\
             & + morph.    & 0.175 & $+0.172$ & 0.032 & 0.87 & 1.13 \\
             & full        & 0.180 & $+0.177$ & 0.031 & 0.88 & 1.13 \\
    \midrule
    BocaSud  & no waves    & 0.023 & $-0.004$ & 0.023 & 0.91 & 1.06 \\
             & waves       & 0.028 & $+0.018$ & 0.022 & 0.92 & 1.06 \\
             & + roughness & 0.025 & $+0.012$ & 0.023 & 0.92 & 1.07 \\
             & + morph.    & 0.023 & $+0.009$ & 0.021 & 0.93 & 1.08 \\
             & full        & 0.022 & $+0.003$ & 0.022 & 0.93 & 1.09 \\
    \midrule
    AltaVila & no waves    & 0.116 & $+0.108$ & 0.042 & 0.81 & 1.12 \\
    Est      & waves       & 0.136 & $+0.130$ & 0.039 & 0.82 & 1.14 \\
             & + roughness & 0.130 & $+0.124$ & 0.039 & 0.82 & 1.16 \\
             & + morph.    & 0.124 & $+0.117$ & 0.042 & 0.81 & 1.20 \\
             & full        & 0.102 & $+0.091$ & 0.045 & 0.79 & 1.25 \\
    \bottomrule
  \end{tabular}
\end{table}

\subsection{Lagrangian transport}
\label{sec:res_drifters}

Thirty-five drifters across twelve deployments were scored against each member
over the interval each drifter was in the water, which ranged from 0.5 to
7.2~hours. Table~\ref{tab:lagrangian} summarises the four metrics and
Figure~\ref{fig:lagrangian} shows their distribution.

Mean Liu--Weisberg skill was 0.539 without wave coupling, 0.505 and 0.506 for
the two wave-coupled fixed-bed members, 0.570 with morphodynamics alone and
0.659 with both morphodynamics and distributed roughness. Endpoint separation
followed the same ordering, from 549 to 583~m for the first four members and
343~m for the last. At drifter level the five distributions overlap heavily
(Figure~\ref{fig:lagrangian}a), so the ordering of the means rests on paired
comparison rather than on separation between groups.

\begin{table}[H]
  \centering
  \caption{Lagrangian metrics by ensemble member, 35 drifters over twelve
  deployments. Path ratio is simulated over observed path length and speed
  ratio is mean speed along the simulated track over the observed
  0.143~\si{\metre\per\second}, both scored on the interval each drifter was in
  the water.}
  \label{tab:lagrangian}
  \begin{tabular}{lrrrr}
    \toprule
    Member & LW skill & Endpoint sep.\ (m) & Path ratio & Speed ratio \\
    \midrule
    No waves            & 0.539 & 549 & 0.78 & 0.82 \\
    Waves               & 0.505 & 569 & 0.80 & 0.84 \\
    Waves + roughness   & 0.506 & 583 & 0.81 & 0.85 \\
    Waves + morph.      & 0.570 & 566 & 0.75 & 0.76 \\
    Full                & \textbf{0.659} & \textbf{343} & \textbf{0.96} & \textbf{0.99} \\
    \bottomrule
  \end{tabular}
\end{table}

The wave penalty is not carried by a few deployments
(Figure~\ref{fig:lagrangian}b). Skill falls with wave coupling in 28 of the 35
drifters (sign test $p = 0.0005$) and in 10 of the 12 deployments (sign test
$p = 0.039$, Wilcoxon on deployment means $p = 0.043$). The two deployments in
which waves help are internally consistent, with all three of their drifters
improving in each, so they represent a genuine difference in behaviour rather
than scatter. The largest single degradation is $-0.133$ and the largest single
improvement $+0.069$.

The observed drifters travelled at a mean speed of 0.143~\si{\metre\per\second}.
Along their own simulated trajectories the members reached 0.118, 0.121, 0.122,
0.109 and 0.142~\si{\metre\per\second} respectively
(Figure~\ref{fig:lagrangian}c), so only the member combining morphodynamics and
distributed roughness reproduces the observed transport rate. The remainder
fall 15 to 24\% short. The same ordering appears in the fraction of drifters
whose path ratio falls below 0.7, which is 31 to 34\% for the first four
members against 11\% for the fifth.

\begin{figure}[H]
  \centering
  \includegraphics[width=\textwidth]{lagrangian_results}
  \caption{Lagrangian validation. (a) Liu--Weisberg skill for every drifter
  under every member, with the horizontal bar giving the mean. (b) The change
  in skill from wave coupling, by deployment, with individual drifters
  overlaid and the ensemble mean dashed. (c) Mean speed along each member's own
  track against the observed value, with the ratio annotated. The member
  combining distributed roughness and an active bed is drawn in colour
  throughout.}
  \label{fig:lagrangian}
\end{figure}

\subsection{Process attribution}
\label{sec:res_attribution}

Figure~\ref{fig:attribution} gives the paired per-drifter contrasts with
bootstrap confidence intervals. Wave coupling reduces Liu--Weisberg skill by
0.035 (95\% CI $-0.053$ to $-0.017$, $p = 0.001$). Distributed roughness on a
fixed bed changes it by 0.001 (CI $-0.016$ to $+0.018$, $p = 0.34$).
Morphodynamics under uniform roughness changes it by 0.065 (CI $-0.019$ to
$+0.150$, $p = 0.23$). The two contrasts in which both treatments are present
are the only large ones with intervals clear of zero. Distributed roughness
with morphodynamics active gives $+0.089$ (CI $+0.046$ to $+0.133$,
$p < 0.001$), and morphodynamics with distributed roughness active gives
$+0.153$ (CI $+0.056$ to $+0.246$, $p = 0.005$).

Neither treatment produces a demonstrable change on its own and their joint
effect exceeds the sum of the individual point estimates, so the two interact
rather than contribute additively.

The joint effect is not evenly distributed across the dataset. It exceeds
$+0.27$ in five deployments, is close to zero in four, and is negative in
three, reaching $-0.39$ in the deployment with the shortest observed track.

\begin{figure}[H]
  \centering
  \includegraphics[width=\textwidth]{ensemble_attribution}
  \caption{Process attribution over the ensemble. (a) and (b) give mean skill
  and endpoint separation against the roughness treatment, one line per bed
  treatment, with the no-wave member as the dashed reference and its bootstrap
  interval shaded. (c) gives paired effect sizes with 95\% bootstrap intervals, where
  filled markers denote intervals clear of zero. (d) gives the per-deployment
  mean of the two headline contrasts.}
  \label{fig:attribution}
\end{figure}

\subsection{Three-dimensional structure}
\label{sec:res_3d}

Over 4043 wet cells in the lagoon interior, with the inlets excluded, mean
current speed increases monotonically from the bed to the surface in every
member (Figure~\ref{fig:profile}). Bed-layer speed ranges from 0.021 to
0.040~\si{\metre\per\second} and surface-layer speed from 0.075 to
0.102~\si{\metre\per\second}, giving surface-to-bed ratios of 1.85, 2.07, 2.05,
4.48 and 3.77 across the five members.

The three fixed-bed members give nearly identical profiles, with ratios between
2.05 and 2.07. Both morphodynamic members roughly double the ratio, and they do
so through the bed layer rather than the surface, with near-bed speed falling
from 0.038 to 0.021~\si{\metre\per\second} while surface speed changes by less.
Adding distributed roughness to an active bed lowers surface speed from 0.102
to 0.078~\si{\metre\per\second}.

Lagoon-mean surface speed does not order the members the way transport does.
The member reproducing the observed drifter speed has a lagoon-mean surface
speed of 0.078~\si{\metre\per\second}, below the 0.080 of a member whose
particles travel 15\% too slowly, and below the 0.102 of the member with the
lowest path ratio of all.

\begin{figure}[H]
  \centering
  \includegraphics[width=0.88\textwidth]{layer_profile_lagoon}
  \caption{Vertical structure in the lagoon interior over the drifter window.
  (a) Mean speed by sigma layer, colour giving the bed treatment and line style
  the roughness treatment. (b) The ratio of surface to bed speed.}
  \label{fig:profile}
\end{figure}

\subsection{Response of the roughness classes}
\label{sec:res_classes}

Cells were grouped by the trachytope class the model assigned them, taking the
dominant class at the nearest roughness link. Comparing classes within a single
member does not isolate the roughness treatment, because the classes occupy
different parts of the basin at different depths. The uniform-roughness members
make this explicit. With no roughness distinction applied at all, their
\emph{Posidonia} cells are still slower at the bed than their sand cells, in a
ratio of 0.86 and 0.87. The roughness effect is therefore read between members
within a class rather than between classes within a member.

On a fixed bed, moving from uniform to distributed roughness reduces mean bed
speed by 1.3\% over sand, 1.1\% over \emph{Cymodocea}, 3.6\% over
\emph{Posidonia} and 4.5\% over rock (Figure~\ref{fig:classes}b). The effect on
the wave field is larger and ordered as canopy attenuation predicts. Mean
orbital velocity at the bed falls by 11.3\% over sand, 13.7\% over
\emph{Cymodocea} and 14.3\% over \emph{Posidonia}, and by only 3.3\% over rock
(Figure~\ref{fig:classes}c), the rock class being the exposed substrate near
the inlets where mean orbital velocity reaches 0.27~\si{\metre\per\second}
against 0.05 to 0.09~\si{\metre\per\second} elsewhere in the basin.

Introducing bed mobility inverts the class ordering of near-bed speed. The
\emph{Posidonia} to sand ratio rises from 0.85 on a fixed bed to 1.10 and 1.27
in the two morphodynamic members, because sand cells lose roughly half their
near-bed speed while \emph{Posidonia} cells lose about a third.

\begin{figure}[H]
  \centering
  \includegraphics[width=\textwidth]{class_speed_uorb}
  \caption{Flow and wave response by trachytope class in the lagoon interior.
  (a) Mean bed and surface speed by class for every member, the full member
  drawn opaque. (b) Change in mean bed speed from uniform to distributed
  roughness on a fixed bed. (c) The same for mean wave orbital velocity at the
  bed.}
  \label{fig:classes}
\end{figure}
